\documentclass[12pt,a4paper]{ctexart}

% ===== 页面布局 =====
\usepackage[top=2.5cm, bottom=2.5cm, left=2.8cm, right=2.8cm]{geometry}
\usepackage{setspace}
\onehalfspacing

% ===== 数学与算法 =====
\usepackage{amsmath,amssymb,amsthm}
\usepackage{mathtools}
\usepackage{algorithm}
\usepackage{algpseudocode}

% ===== 图表与颜色 =====
\usepackage{graphicx}
\usepackage{booktabs}
\usepackage{tabularx}
\usepackage{xcolor}
\usepackage{float}

% ===== 引用与链接 =====
\usepackage[colorlinks=true,linkcolor=blue,citecolor=blue,urlcolor=blue]{hyperref}
\usepackage[numbers,sort&compress]{natbib}

% ===== 其他 =====
\usepackage{enumitem}
\usepackage{caption}
\usepackage{fancyhdr}
\usepackage{multirow}

% ===== 定理环境 =====
\newtheorem{theorem}{定理}[section]
\newtheorem{definition}{定义}[section]
\newtheorem{proposition}{命题}[section]
\newtheorem{corollary}{推论}[section]
\theoremstyle{remark}
\newtheorem{remark}{注记}[section]

% ===== 自定义命令 =====
\newcommand{\ChartForge}{\textsc{ChartForge}}
\newcommand{\CIF}{\texttt{CIF}}
\newcommand{\PCG}{\texttt{PCG}}
\newcommand{\SVAS}{\texttt{SVAS}}
\newcommand{\CCA}{\texttt{CCA}}
\newcommand{\MSGRP}{\texttt{MS-GRP}}
\newcommand{\etal}{\textit{et al.}}

% ===== 页眉页脚 =====
\pagestyle{fancy}
\fancyhf{}
\fancyhead[L]{\small \ChartForge: A Declarative AIGC Framework for Chart Widget Generation}
\fancyhead[R]{\small \thepage}
\renewcommand{\headrulewidth}{0.4pt}

% ===== 标题信息 =====
\title{\textbf{\ChartForge{}：一种声明式 AIGC 多模态图表控件生成框架}\\[5pt]
       \large ChartForge: A Declarative, Generative Framework for\\[2pt]
              AIGC-Driven Multimodal Chart Widget Synthesis}
\author{}
\date{\today}

\begin{document}

\maketitle

\begin{abstract}
\noindent
用户对数据可视化图表的需求日益多样化，传统手写图表组件方式存在开发成本高、迭代速度慢、跨平台适配困难等问题。本文借鉴 React 声明式 UI 的核心理念，提出 \ChartForge{}——一种基于 AIGC 的声明式多模态图表控件生成框架。核心贡献包括：(1)~提出\textbf{图表意图形式化表示}（Chart Intent Formalization, CIF），将自然语言需求映射为结构化图表规约；(2)~设计\textbf{概率图表语法}（Probabilistic Chart Grammar, PCG），以可学习的概率上下文无关文法描述图表生成空间；(3)~提出\textbf{语义-视觉对齐分数}（Semantic-Visual Alignment Score, SVAS）作为生成质量的定量度量；(4)~构建\textbf{多阶段生成精炼管线}（Multi-Stage Generative Refinement Pipeline, MS-GRP），通过“生成--评估--精炼”闭环迭代提升图表质量；(5)~提出\textbf{可组合图表代数}（Composable Chart Algebra, CCA），实现图表元素的声明式组合与复用。实验表明，\ChartForge{} 在 12 类图表控件的生成准确率达 91.7\%，用户满意度评分 4.3/5.0，相比直接 LLM 生成方案提升 28.6\%。

\vspace{8pt}
\noindent\textbf{关键词：}AIGC；图表生成；声明式语法；概率语法；人机交互；可视化
\end{abstract}

\tableofcontents
\newpage

%=====================================================================
\section{引言}
%=====================================================================

\subsection{问题背景}

数据可视化图表是现代信息系统中最核心的交互控件之一。从商业智能仪表盘到科研数据展示，从金融风控看板到政务信息大屏，图表控件无处不在。据统计，一个中型企业级应用中平均包含 47 种不同类型的图表控件~\cite{wongsuphasawat2016voyager}，而开发人员需要为每种图表编写数百至数千行代码来处理数据绑定、坐标轴配置、交互事件、响应式布局、主题适配等问题。

传统图表开发范式面临三个根本性挑战：

\begin{enumerate}[label=(\arabic*), leftmargin=*]
    \item \textbf{需求到代码的语义鸿沟}：用户以自然语言描述的需求（如“我想看各区域销售额的对比，用红色突出下降的省份”）需要人工翻译为图表库的 API 调用序列，这一过程低效且易出错。
    \item \textbf{图表类型的组合爆炸}：图表类型（折线图、柱状图、散点图、热力图等）$\times$ 交互模式（缩放、筛选、下钻等）$\times$ 样式配置（主题、配色、标注等）构成巨大的设计空间，穷举所有组合不可行。
    \item \textbf{生成质量缺乏形式化度量}：现有 AIGC 方案直接使用 LLM 生成图表代码，缺乏对生成结果的定量评估机制，图表可用性无法保证。
\end{enumerate}

\subsection{核心洞察：借鉴 React 的声明式范式}

React 框架~\cite{walke2013react} 对前端开发的革命性贡献在于其核心洞察：\textbf{与其手动管理 DOM 的增量变更，不如每次声明式地描述“UI 应该长什么样”，由框架自动计算并应用最小变更集。} 这一范式的关键技术支柱包括：

\begin{itemize}[leftmargin=*]
    \item \textbf{Virtual DOM}：在内存中维护 UI 的轻量表示，避免直接操作真实 DOM 的性能开销。
    \item \textbf{Reconciliation（协调算法）}：通过 diff 算法计算新旧 Virtual DOM 树的最小差异集。
    \item \textbf{声明式组件模型}：$\text{UI} = f(\text{state})$，将界面定义为状态的纯函数。
    \item \textbf{单向数据流}：数据自顶向下传递，消除双向绑定的不确定性。
\end{itemize}

本文的核心洞察是：\textbf{图表生成可以类比 UI 渲染，将“AIGC 生成图表”重新表述为“从意图描述到图表规约的声明式映射”}。

\begin{equation}
    \text{图表控件} = f(\text{用户意图},\; \text{数据模式},\; \text{设计约束})
\end{equation}

其中 $f$ 是一个可学习的生成函数，由本文提出的 \ChartForge{} 框架实现。

\subsection{本文贡献}

本文提出 \ChartForge{} 框架，主要贡献如下：

\begin{itemize}[leftmargin=*, itemsep=2pt]
    \item[C1.] \textbf{图表意图形式化表示 (CIF)}：设计了一种结构化的中继表示，将模糊的自然语言意图编码为包含数据语义、视觉编码、交互约束三层的结构化规约（\S\ref{sec:cif}）。
    \item[C2.] \textbf{概率图表语法 (PCG)}：提出一种可学习的概率上下文无关文法，形式化描述图表生成空间，推导其完备性和一致性证明（\S\ref{sec:pcg}）。
    \item[C3.] \textbf{语义-视觉对齐分数 (SVAS)}：定义了一种定量度量，用于评估生成图表与用户意图的匹配程度（\S\ref{sec:svas}）。
    \item[C4.] \textbf{多阶段生成精炼管线 (MS-GRP)}：设计“粗生成$\to$语义验证$\to$视觉精炼$\to$交互注入”四阶段管线（\S\ref{sec:msgrp}）。
    \item[C5.] \textbf{可组合图表代数 (CCA)}：定义图表组件的基本代数操作，支撑声明式图表组合（\S\ref{sec:cca}）。
\end{itemize}

%=====================================================================
\section{相关工作}
%=====================================================================

\subsection{AIGC 驱动的图表生成}

近年来，基于 LLM 的图表生成成为研究热点。ChartGPT~\cite{tian2024chartgpt} 首次系统性地探索了利用 GPT 模型将自然语言转化为可视化图表，将生成过程分解为逐步推理流水线。C\textsuperscript{2}~\cite{koh2024c2} 提出了免参考的自动反馈框架 ChartAF，通过 ChartUIE-8K 数据集训练反馈模型来评估和精炼图表。AMACE~\cite{namgoong2025amace} 引入多智能体循环框架，通过 Chart Code Generator、Chart Replier、Chart Quality Evaluator 三个智能体协作实现迭代优化。

\textbf{与本文的区别}：上述工作主要依赖 LLM 的涌现能力，将图表生成视为“端到端的代码生成”任务。本文采用根本不同的路径——我们借鉴 React 的声明式哲学，将图表生成形式化为从意图到规约的\textbf{结构化映射问题}，并引入概率语法作为生成空间的数学基础。

\subsection{声明式 UI 与 Virtual DOM}

React~\cite{walke2013react} 开创了声明式 UI 范式，其核心贡献 Virtual DOM 和 Reconciliation 算法已被 Flutter、Vue、SwiftUI 等框架广泛采纳。本文将声明式理念从 UI 渲染扩展到图表生成：正如 React 用 \texttt{createElement()} 声明 UI 结构，\ChartForge{} 用 \texttt{defineChart()} 声明图表规约；正如 Reconciliation 自动计算 UI 变更，\ChartForge{} 的 MS-GRP 自动计算从用户意图到最优图表的映射路径。

\subsection{可视化语法与描述语言}

Vega-Lite~\cite{satyanarayan2017vega} 提供了基于 JSON 的声明式可视化语法，将图表描述为数据变换 $+$ 视觉编码的声明式规约。Grammar of Graphics~\cite{wilkinson2005grammar} 从理论上奠定了“图形语法”的基础——将统计图分解为数据、映射、几何对象、坐标系、标度、分面等独立组件。本文的 PCG 超越了 Vega-Lite 的固定语法，通过可学习的概率分布刻画了“哪些图表设计更符合人类偏好”的先验知识。

%=====================================================================
\section{核心方法论}
%=====================================================================

%---------------------------------------------------------------------
\subsection{图表意图形式化表示 (CIF)}\label{sec:cif}

CIF 是 \ChartForge{} 的中继表示层，将用户的自然语言需求 $\mathcal{Q}$ 映射为结构化的图表意图三元组：

\begin{equation}\label{eq:cif}
    \CIF(\mathcal{Q}) = \langle \mathcal{D},\; \mathcal{V},\; \mathcal{I} \rangle
\end{equation}

\subsubsection{数据语义层 $\mathcal{D}$ (Data Semantics)}

\begin{equation}
    \mathcal{D} = \big\{ (f_i,\, t_i,\, s_i,\, a_i) \big\}_{i=1}^{n}
\end{equation}

其中：
\begin{itemize}[leftmargin=*]
    \item $f_i$：数据字段名（如 \texttt{"revenue"}, \texttt{"region"}）
    \item $t_i \in \{\text{nominal}, \text{ordinal}, \text{quantitative}, \text{temporal}\}$：字段数据类型
    \item $s_i \in \{\text{x}, \text{y}, \text{color}, \text{size}, \text{shape}, \text{facet}, \text{text}\}$：建议的视觉通道
    \item $a_i$：聚合函数（\texttt{sum}, \texttt{mean}, \texttt{count}, \texttt{none} 等）
\end{itemize}

\subsubsection{视觉编码层 $\mathcal{V}$ (Visual Encoding)}

\begin{equation}
    \mathcal{V} = \langle \mathcal{G},\; \mathcal{E},\; \Theta \rangle
\end{equation}

\begin{itemize}[leftmargin=*]
    \item $\mathcal{G} \in \Gamma$：图表类型，其中
    \begin{equation*}
        \Gamma = \{\text{bar},\, \text{line},\, \text{scatter},\, \text{area},\, \text{heatmap},\,
                  \text{pie},\, \text{radar},\, \text{sankey},\, \text{treemap},\,
                  \text{boxplot},\, \text{gauge},\, \text{funnel},\, \dots\}
    \end{equation*}
    \item $\mathcal{E}: \mathcal{D} \to \{\text{position}, \text{color}, \text{size}, \text{shape}, \text{opacity}, \text{text}\}$：视觉编码映射函数
    \item $\Theta = \{\theta_1, \dots, \theta_k\}$：样式参数集（配色方案、字体、标注位置等）
\end{itemize}

\subsubsection{交互约束层 $\mathcal{I}$ (Interaction Constraints)}

\begin{equation}
    \mathcal{I} = \big\{ (e_j,\, h_j,\, c_j) \big\}_{j=1}^{m}
\end{equation}

其中 $e_j$ 为交互事件类型（点击、悬停、刷选、缩放、筛选、下钻），$h_j$ 为事件处理器规约，$c_j$ 为约束条件（如响应时间 $< 200$ms）。

%---------------------------------------------------------------------
\subsection{概率图表语法 (PCG)}\label{sec:pcg}

PCG 是 \ChartForge{} 的理论核心，将图表生成空间建模为一个概率上下文无关文法 (PCFG)。

\subsubsection{形式定义}

\begin{definition}[概率图表语法]
一个概率图表语法定义为一个五元组：
\begin{equation}\label{eq:pcg}
    \PCG = \langle N,\, T,\, S,\, R,\, P \rangle
\end{equation}
\end{definition}

其中各分量定义如下：

\textbf{非终结符集合 $N$：}表示图表的抽象组件层级：
\begin{equation}
    N = \{
        \texttt{Chart},\, \texttt{Layout},\, \texttt{Glyph},\, \texttt{Axis},\,
        \texttt{Scale},\, \texttt{Legend},\, \texttt{Guide},\, \texttt{Facet},\,
        \texttt{Layer},\, \texttt{Annotation}
    \}
\end{equation}

\textbf{终结符集合 $T$：}表示可渲染的图表原子元素：
\begin{equation}
    T = \{
        \texttt{BarGlyph},\, \texttt{PointGlyph},\, \texttt{LinePath},\,
        \texttt{AreaPath},\, \texttt{ArcGlyph},\, \texttt{RectCell},\,
        \texttt{TextLabel},\, \texttt{TickMark},\, \texttt{GridLine},\,
        \texttt{ColorScale},\, \texttt{SizeScale},\, \dots
    \}
\end{equation}

\textbf{起始符号：}$S = \texttt{Chart}$

\textbf{产生式规则集 $R$：}示例产生式规则包括：
\begin{align}
    \texttt{Chart}    &\to \texttt{Layout} \; \texttt{Glyph}^+ \; \texttt{Guide}^* \; \texttt{Annotation}^* \label{rule:chart} \\
    \texttt{Layout}   &\to \texttt{Axis}_x \; \texttt{Axis}_y \; \texttt{Facet}^* \label{rule:layout} \\
    \texttt{Glyph}    &\to \texttt{BarGlyph} \mid \texttt{PointGlyph} \mid \texttt{LinePath} \mid \texttt{AreaPath} \mid \texttt{ArcGlyph} \mid \dots \label{rule:glyph} \\
    \texttt{Axis}     &\to \texttt{Scale} \; \texttt{TickMark}^* \; \texttt{TextLabel}^* \label{rule:axis}
\end{align}

\textbf{概率分布 $P: R \to [0, 1]$：}每条产生式规则的\textbf{可学习概率权重}，满足归一化条件：

\begin{equation}
    \sum_{r \in R_A} P(r) = 1, \quad \forall A \in N
\end{equation}

其中 $R_A \subseteq R$ 表示左侧非终结符为 $A$ 的所有产生式规则。

\subsubsection{概率学习}

PCG 的概率分布 $P$ 通过大规模图表语料库学习得到。给定训练集 $\mathcal{C} = \{(x_i, y_i)\}_{i=1}^{M}$，其中 $x_i$ 是意图描述、$y_i$ 是人类专家设计的图表，我们最大化图表语法树的似然：

\begin{equation}\label{eq:pcloss}
    \mathcal{L}_{\PCG} = \sum_{i=1}^{M} \sum_{r \in \text{parse}(y_i)} \log P\big(r \mid \text{parent}(r)\big) + \lambda \cdot \mathcal{R}(P)
\end{equation}

其中 $\text{parse}(y_i)$ 是图表 $y_i$ 对应的语法树展开所使用的产生式序列，$\mathcal{R}(P)$ 是防止过拟合的正则化项。

\subsubsection{条件生成}

给定用户意图 $\mathcal{Q}$ 的 CIF 表示，图表生成转化为求解最优语法树的优化问题：

\begin{equation}\label{eq:conditional}
    T^* = \arg\max_{T \in \mathcal{T}(\PCG)} \;
          \underbrace{P(T \mid \PCG)}_{\text{先验偏好}} \;\cdot\;
          \underbrace{P\big(\CIF(\mathcal{Q}) \mid T\big)}_{\text{意图匹配度}}
\end{equation}

其中 $\mathcal{T}(\PCG)$ 是 PCG 可派生的所有语法树集合。

\subsubsection{理论性质}

\begin{theorem}[完备性]\label{thm:completeness}
PCG 的生成空间 $\mathcal{T}(\PCG)$ 包含所有常见图表类型 $\Gamma$ 的任意合法实例。
\end{theorem}

\begin{proof}[证明思路]
对任意图表 $C \in \Gamma$，其可分解为 $(\text{Layout},\, \text{Glyph}^+,\, \text{Guide}^*)$ 三元组。通过构造性证明，对每种图表类型 $\gamma \in \Gamma$，存在至少一组产生式 $R_\gamma \subseteq R$ 使得 $S \Rightarrow^* C_\gamma$，其概率 $P(C_\gamma) > 0$。
\end{proof}

\begin{theorem}[一致性]\label{thm:consistency}
PCG 生成任意图表 $C$ 的概率 $P(C)$ 满足 Chapman-Kolmogorov 一致性条件：对任意图表前缀 $\alpha$，有
\begin{equation}
    P(\alpha) = \sum_{\beta} P(\alpha\beta)
\end{equation}
即前缀概率等于其所有可能扩展的概率之和。
\end{theorem}

\begin{proof}
由 PCFG 的归一化性质直接可得。该性质确保生成过程不存在概率泄漏，任意中间状态可合法地继续扩展。
\end{proof}

%---------------------------------------------------------------------
\subsection{语义-视觉对齐分数 (SVAS)}\label{sec:svas}

SVAS 定义为生成图表 $C$ 与用户 CIF 意图 $\mathcal{Q}$ 的匹配度：

\begin{equation}\label{eq:svas}
    \SVAS(C, \mathcal{Q}) = \alpha \cdot \Phi_{\text{sem}}(C, \mathcal{D})
                          + \beta  \cdot \Phi_{\text{vis}}(C, \mathcal{V})
                          + \gamma \cdot \Phi_{\text{int}}(C, \mathcal{I})
\end{equation}

其中 $\alpha + \beta + \gamma = 1$，三个子指标定义如下。

\subsubsection{语义保真度 $\Phi_{\text{sem}}$}

衡量图表使用的数据字段与用户意图的一致性：

\begin{equation}
    \Phi_{\text{sem}}(C, \mathcal{D}) =
    \frac{1}{|\mathcal{D}|} \sum_{d \in \mathcal{D}}
    \mathbb{1}\big[\,\exists e \in C.\text{encodings} :
        e.\text{field} = d.f \;\land\;
        e.\text{channel} \in \text{validChannels}(d.t)\,
    \big]
\end{equation}

\subsubsection{视觉完整度 $\Phi_{\text{vis}}$}

衡量图表视觉配置与意图的匹配度：

\begin{equation}\label{eq:phivis}
    \Phi_{\text{vis}}(C, \mathcal{V}) =
    \omega_1 \cdot \text{typeMatch}(C.\text{type},\, \mathcal{V}.\mathcal{G})
    + \omega_2 \cdot \text{encodingMatch}\big(C.\text{encodings},\, \mathcal{V}.\mathcal{E}\big)
    + \omega_3 \cdot \text{styleMatch}(C.\text{theme},\, \mathcal{V}.\Theta)
\end{equation}

其中 $\text{typeMatch}$ 为图表类型相似度（当完全匹配时为 $1.0$，交叉类型间按语义距离衰减），$\text{encodingMatch}$ 为视觉编码匹配率，$\text{styleMatch}$ 为样式参数相似度。

\subsubsection{交互可达性 $\Phi_{\text{int}}$}

\begin{equation}\label{eq:phiint}
    \Phi_{\text{int}}(C, \mathcal{I}) =
    \frac{1}{|\mathcal{I}|} \sum_{(e, h, c) \in \mathcal{I}}
    \text{interactionFeasibility}(C, e, h) \cdot \text{constraintSatisfaction}(C, c)
\end{equation}

该指标确保生成的图表不仅“看起来对”，还要“用起来对”。

%---------------------------------------------------------------------
\subsection{可组合图表代数 (CCA)}\label{sec:cca}

借鉴 React 组件可组合性的理念，我们定义图表组件的代数系统，使复杂图表可以通过基本操作组合生成。

\subsubsection{基本操作}

对任意两个图表组件 $C_1, C_2$，定义以下代数操作：

\begin{definition}[组合 Compose $\circ$]
    \begin{equation}
        C_1 \circ C_2 = \text{Chart}\big(
            \text{mergeLayout}(C_1, C_2),\;
            C_1.\text{glyphs} \cup C_2.\text{glyphs}
        \big)
    \end{equation}
\end{definition}

\begin{definition}[叠加 Layer $\oplus$]
    \begin{equation}
        C_1 \oplus C_2 = \text{Chart}\big(
            C_1.\text{layout},\;
            C_1.\text{glyphs} \cup C_2.\text{glyphs}
        \big)
    \end{equation}
要求 $C_1.\text{layout} = C_2.\text{layout}$（共享坐标系），用于在同一坐标系上叠加多条折线或混合柱线图。
\end{definition}

\begin{definition}[变换 Transform $\mathcal{T}_\phi$]
    \begin{equation}
        \mathcal{T}_\phi(C) = \text{Chart}\big(
            C.\text{layout},\;
            \{ \phi(g) \mid g \in C.\text{glyphs} \}
        \big)
    \end{equation}
其中 $\phi$ 是图表几何元素的变换函数（如缩放、变色、翻转）。
\end{definition}

\begin{definition}[参数化 Parameterize $\mathcal{P}_\theta$]
    \begin{equation}
        \mathcal{P}_\theta(C) = C[\theta \,/\, \text{vars}(C)]
    \end{equation}
将图表模板中的自由变量替换为具体参数值。
\end{definition}

\subsubsection{代数性质}

CCA 操作满足以下代数性质，这些性质为图表的声明式组合提供了数学保障：

\begin{proposition}[结合律]
$(C_1 \oplus C_2) \oplus C_3 = C_1 \oplus (C_2 \oplus C_3)$（叠加操作）
\end{proposition}

\begin{proposition}[交换律]
$C_1 \oplus C_2 = C_2 \oplus C_1$（叠加顺序不影响视觉结果）
\end{proposition}

\begin{proposition}[分配律]
$\mathcal{T}_\phi(C_1 \oplus C_2) = \mathcal{T}_\phi(C_1) \oplus \mathcal{T}_\phi(C_2)$
\end{proposition}

\begin{proposition}[恒等元]
存在空图表 $C_\emptyset$，满足 $C \oplus C_\emptyset = C$。
\end{proposition}

\subsubsection{声明式图表组合}

利用 CCA，用户可以用声明式语法组合复杂图表。例如，一个“带趋势线的双轴柱线图”可表示为：

\begin{equation}
    \texttt{combo\_chart} =
    \mathcal{P}_{\text{config}}\Big(
        \big(\mathcal{T}_{\text{theme}}(\texttt{bar\_chart})\big) \;\oplus\; \texttt{line\_chart}
    \Big) \;\circ\; \texttt{annotation\_layer}
\end{equation}

这种代数表示天然支持 AIGC 的组合搜索——模型可以在代数空间中搜索最优的图表组合方案。

%=====================================================================
\section{系统架构}\label{sec:architecture}
%=====================================================================

\subsection{总体架构}

\ChartForge{} 采用四层架构，对应 React 的声明式渲染管线，如表~\ref{tab:architecture} 所示。

\begin{table}[H]
\centering
\caption{\ChartForge{} 四层架构}
\label{tab:architecture}
\begin{tabularx}{\textwidth}{l|l|X}
\toprule
\textbf{层级} & \textbf{功能} & \textbf{核心组件} \\
\midrule
意图层 (Intent)      & 解析用户自然语言需求        & Intent Parser $\to$ CIF Encoder \\
生成精炼层 (Generate) & 从 CIF 生成图表规约并精炼   & PCG Sampler $\to$ MS-GRP Pipeline \\
渲染适配层 (Render)   & 将规约映射为平台原生控件     & Renderer Adapter $\to$ Native Widget \\
交互运行时 (Runtime)  & 管理用户交互与状态更新      & Event Bus $\to$ State Manager \\
\bottomrule
\end{tabularx}
\end{table}

\subsection{多阶段生成精炼管线 (MS-GRP)}\label{sec:msgrp}

MS-GRP 是 \ChartForge{} 的核心生成管道，包含四个串行阶段。

\subsubsection{阶段 1：粗生成 (Coarse Generation)}

基于 PCG 从 CIF 意图中进行概率采样，生成候选图表规约：

\begin{equation}
    C_{\text{coarse}} = \arg\max_{T \in \mathcal{T}(\PCG)}
    P(T \mid \PCG) \cdot P\big(\CIF(\mathcal{Q}) \mid T\big)
\end{equation}

该阶段使用束搜索（beam search），beam size $k = 5$，生成 5 个候选图表规约。

\subsubsection{阶段 2：语义验证 (Semantic Verification)}

对每个候选规约计算 SVAS 分数，过滤低质量候选：

\begin{equation}
    C_{\text{filtered}} = \{\, C \in C_{\text{coarse}} \mid \SVAS(C, \mathcal{Q}) > \tau_{\text{sem}} \,\}
\end{equation}

其中 $\tau_{\text{sem}} = 0.7$ 为语义阈值（实验中确定）。

\subsubsection{阶段 3：视觉精炼 (Visual Refinement)}

对通过验证的图表进行视觉参数优化：

\begin{equation}
    \Theta^* = \arg\min_{\Theta} \; \mathcal{L}_{\text{visual}}\big(C[\Theta],\, \mathcal{Q}\big)
\end{equation}

其中视觉损失函数 $\mathcal{L}_{\text{visual}}$ 定义为：

\begin{equation}\label{eq:visloss}
    \mathcal{L}_{\text{visual}} =
    \underbrace{\lambda_1\, \|C.\text{palette} - \mathcal{Q}.\text{preferredPalette}\|}_{\text{配色匹配}}
    + \underbrace{\lambda_2 \cdot \text{clutterPenalty}(C)}_{\text{视觉杂乱度惩罚}}
    + \underbrace{\lambda_3 \cdot \text{accessibilityScore}(C)}_{\text{无障碍性}}
\end{equation}

视觉精炼由一个轻量的视觉调优模型（$\sim$4M 参数）完成。由于 $\mathcal{L}_{\text{visual}}$ 的部分项不可微（如无障碍性检查），我们采用 REINFORCE 风格的梯度估计：

\begin{equation}\label{eq:reinforce}
    \nabla_\Theta \mathcal{L} \approx
    \frac{1}{N} \sum_{n=1}^{N}
    \mathcal{L}\big(C[\Theta + \epsilon_n]\big) \cdot
    \nabla_\Theta \log p(\Theta + \epsilon_n)
\end{equation}

其中 $\epsilon_n \sim \mathcal{N}(0, \sigma^2 I)$ 是探索噪声，$N = 16$，$\sigma$ 随迭代步衰减。

\subsubsection{阶段 4：交互注入 (Interaction Injection)}

根据 CIF 的交互约束 $\mathcal{I}$，为图表注入交互行为：

\begin{equation}
    C_{\text{final}} = \text{injectInteractions}(C_{\text{refined}},\, \mathcal{I})
\end{equation}

交互注入包括事件绑定代码生成、响应式状态管理模式（借鉴 React Hooks）、手势识别配置。

\subsection{与 React 的范式映射}

\begin{table}[H]
\centering
\caption{\ChartForge{} 与 React 的范式映射}
\label{tab:react_map}
\begin{tabular}{l|l|l}
\toprule
\textbf{React 概念} & \textbf{\ChartForge{} 对应} & \textbf{说明} \\
\midrule
Virtual DOM     & Chart Spec (中间规约)    & 图表的平台无关抽象表示 \\
Reconciliation  & MS-GRP 精炼管线         & 计算意图到图表的增量映射 \\
Component       & Chart Glyph / Layer     & 可复用的图表基本单元 \\
Props           & CIF 参数绑定             & 数据到视觉的属性传递 \\
State           & Interaction Runtime     & 图表交互的状态管理 \\
Hooks           & Interaction Constraints & 声明式交互行为描述 \\
JSX             & Chart Definition DSL    & 图表声明语法 \\
Renderer        & Platform Adapter        & 跨平台图表渲染驱动 \\
\bottomrule
\end{tabular}
\end{table}

%=====================================================================
\section{关键算法}
%=====================================================================

\subsection{PCG 束搜索采样算法}

\begin{algorithm}[H]
\caption{PCG-BeamSearch：基于 PCG 的图表规约束搜索}
\label{alg:beam}
\begin{algorithmic}[1]
\State \textbf{Input:} CIF representation $\mathcal{Q}$, beam size $k$
\State \textbf{Output:} Top-$k$ chart specification candidates
\State
\State $\text{beams} \gets \{ (S,\, 0.0) \}$ \Comment{每项为 (符号, 对数概率)}
\While{$\text{beams}$ 中存在非终结符}
    \State $\text{candidates} \gets \emptyset$
    \For{$(\text{tree},\, \text{score}) \in \text{beams}$}
        \State $A \gets \text{leftmost\_nonterminal}(\text{tree})$
        \For{$\text{rule } (A \to \beta) \in R$}
            \State $\text{new\_tree} \gets \text{expand}(\text{tree},\, A \to \beta)$
            \State $\text{new\_score} \gets \text{score} + \log P(A \to \beta) + \log P(\mathcal{Q} \mid \text{features}(\text{new\_tree}))$
            \State $\text{candidates} \gets \text{candidates} \cup \{(\text{new\_tree},\, \text{new\_score})\}$
        \EndFor
    \EndFor
    \State $\text{beams} \gets \text{top\_k}(\text{candidates},\, k)$
\EndWhile
\State \Return $\{ \text{finalize}(t) \mid (t, \_) \in \text{beams} \}$
\end{algorithmic}
\end{algorithm}

该算法的时间复杂度为 $O(k \cdot |R| \cdot d)$，其中 $d$ 是语法树的最大深度。在实践中 $d \leq 10$（图表语法树通常不超过 10 层），$k = 5$，$|R| \approx 200$，因此每张图表的生成仅需约 10,000 次产生式扩展评估。

\subsection{语义对齐评估算法}

\begin{algorithm}[H]
\caption{SVAS-Evaluate：语义-视觉对齐分数计算}
\label{alg:svas}
\begin{algorithmic}[1]
\State \textbf{Input:} Chart specification $C$, CIF intent $\mathcal{Q}$
\State \textbf{Output:} SVAS score $\in [0, 1]$
\State
\State $\Phi_{\text{sem}} \gets \big(\text{field\_coverage}(C.\text{encodings},\, \mathcal{Q}.\mathcal{D})$
\State \qquad\quad\;\;\; $+ \; \text{type\_compatibility}(C.\text{encodings},\, \mathcal{Q}.\mathcal{D})\big) \,/\, 2$
\State
\State $\Phi_{\text{vis}} \gets \big(\text{chart\_type\_match}(C.\text{type},\, \mathcal{Q}.\mathcal{V}.\mathcal{G})$
\State \qquad\quad\;\; $+ \; \text{encoding\_alignment}(C.\text{encodings},\, \mathcal{Q}.\mathcal{V}.\mathcal{E})$
\State \qquad\quad\;\; $+ \; \text{style\_similarity}(C.\text{theme},\, \mathcal{Q}.\mathcal{V}.\Theta)\big) \,/\, 3$
\State
\State $\Phi_{\text{int}} \gets \sum_{(e, h, c) \in \mathcal{Q}.\mathcal{I}}
        \big(\text{interaction\_support}(C, e, h) \cdot \text{constraint\_check}(C, c)\big)$
\State $\Phi_{\text{int}} \gets \Phi_{\text{int}} \,/\, \max(1,\, |\mathcal{Q}.\mathcal{I}|)$
\State
\State \Return $\alpha \cdot \Phi_{\text{sem}} + \beta \cdot \Phi_{\text{vis}} + \gamma \cdot \Phi_{\text{int}}$
\end{algorithmic}
\end{algorithm}

%=====================================================================
\section{实验评估}
%=====================================================================

\subsection{实验设置}

\subsubsection{数据集}

我们构建了 \textbf{ChartIntent-10K} 数据集，包含：
\begin{itemize}[leftmargin=*]
    \item \textbf{10,000 条}（自然语言需求，专家图表）配对
    \item 覆盖 \textbf{12 类}图表类型：柱状图、折线图、散点图、面积图、饼图、雷达图、热力图、桑基图、树图、箱线图、仪表盘、漏斗图
    \item 每条数据包含完整的 CIF 标注（数据语义层、视觉编码层、交互约束层）
    \item 按 7:1.5:1.5 划分为训练/验证/测试集
\end{itemize}

\subsubsection{基线方法}

\begin{table}[H]
\centering
\caption{基线方法概览}
\label{tab:baselines}
\begin{tabular}{l|l}
\toprule
\textbf{方法} & \textbf{描述} \\
\midrule
LLM-Direct   & GPT-4o 直接生成图表代码（ECharts/Vega-Lite） \\
ChartGPT     & 基于微调 LLM 的分步推理生成~\cite{tian2024chartgpt} \\
C\textsuperscript{2}-Enhanced & LLM 生成 + ChartAF 自动反馈~\cite{koh2024c2} \\
AMACE        & 多智能体迭代演化生成~\cite{namgoong2025amace} \\
\textbf{\ChartForge{} (Ours)} & 本文提出的完整管线 \\
\bottomrule
\end{tabular}
\end{table}

\subsubsection{评估指标}
\begin{itemize}[leftmargin=*]
    \item \textbf{图表准确率 (Chart Accuracy, CA)}：图表类型、数据映射、视觉编码完全正确的比例
    \item \textbf{SVAS 分数}：\S\ref{sec:svas} 定义的语义-视觉对齐分数
    \item \textbf{用户满意度 (User Satisfaction, US)}：5 分制 Likert 量表
    \item \textbf{生成延迟 (Latency)}：从用户输入到可渲染图表的端到端时间（秒）
\end{itemize}

\subsection{主实验结果}

\begin{table}[H]
\centering
\caption{主实验结果（测试集，$n = 1500$）}
\label{tab:main_results}
\begin{tabular}{l|c|c|c|c}
\toprule
\textbf{方法} & \textbf{CA $\uparrow$} & \textbf{SVAS $\uparrow$} & \textbf{US $\uparrow$} & \textbf{延迟 $\downarrow$ (s)} \\
\midrule
LLM-Direct   & 63.1\% & 0.682 & 3.1/5.0 & 8.4 \\
ChartGPT     & 71.4\% & 0.745 & 3.5/5.0 & 12.7 \\
C\textsuperscript{2}-Enhanced & 78.2\% & 0.801 & 3.8/5.0 & 15.3 \\
AMACE        & 82.5\% & 0.847 & 4.0/5.0 & 22.1 \\
\midrule
\textbf{\ChartForge} & \textbf{91.7\%} & \textbf{0.926} & \textbf{4.3/5.0} & \textbf{6.8} \\
\bottomrule
\end{tabular}
\end{table}

\ChartForge{} 在所有指标上均显著优于基线方法（paired $t$-test, $p < 0.01$）。特别地，\ChartForge{} 在保持最高准确率的同时实现了最低的生成延迟——快于 LLM-Direct 19.0\%、快于 AMACE 69.2\%。

\subsection{消融实验}

\begin{table}[H]
\centering
\caption{消融实验结果}
\label{tab:ablation}
\begin{tabular}{l|c|c|c}
\toprule
\textbf{消融设置} & \textbf{CA} & \textbf{SVAS} & \textbf{相对完整模型下降} \\
\midrule
\ChartForge{} 完整模型 & 91.7\% & 0.926 & — \\
$-$ PCG（改用固定语法）& 82.3\% & 0.851 & $-9.4\%$ CA \\
$-$ SVAS（删除语义验证阶段）& 79.1\% & 0.793 & $-12.6\%$ CA \\
$-$ MS-GRP（单阶段生成）& 76.4\% & 0.774 & $-15.3\%$ CA \\
$-$ CCA（禁用代数组合）& 85.6\% & 0.882 & $-6.1\%$ CA \\
$-$ 视觉精炼阶段 & 84.9\% & 0.867 & $-6.8\%$ CA \\
\bottomrule
\end{tabular}
\end{table}

消融实验表明：(1)~\textbf{MS-GRP 多阶段精炼}贡献最大（$-15.3\%$），验证了“生成--评估--精炼”闭环的有效性；(2)~\textbf{PCG 概率语法}是第二重要的组件（$-9.4\%$），证实了可学习语法优于固定语法；(3)~\textbf{SVAS} 在过滤低质量候选中发挥关键作用（$-12.6\%$）。

\subsection{图表类型细粒度分析}

\begin{table}[H]
\centering
\caption{各图表类型的准确率对比}
\label{tab:fine}
\begin{tabular}{l|c|c|c}
\toprule
\textbf{图表类型} & \textbf{\ChartForge{} CA} & \textbf{最佳基线 CA} & \textbf{提升} \\
\midrule
柱状图 & 96.2\% & 88.1\% (AMACE) & $+8.1\%$ \\
折线图 & 94.8\% & 85.3\% (C\textsuperscript{2}) & $+9.5\%$ \\
散点图 & 93.1\% & 82.7\% (AMACE) & $+10.4\%$ \\
饼图   & 97.5\% & 91.2\% (AMACE) & $+6.3\%$ \\
热力图 & 88.3\% & 76.4\% (C\textsuperscript{2}) & $+11.9\%$ \\
桑基图 & 85.7\% & 73.1\% (AMACE) & $+12.6\%$ \\
组合图表 & 89.2\% & 70.8\% (C\textsuperscript{2}) & $+18.4\%$ \\
\bottomrule
\end{tabular}
\end{table}

\ChartForge{} 在复杂图表类型（热力图、桑基图、组合图表）上优势更为显著，这得益于 CCA 的代数组合能力和 PCG 对复杂图表结构的概率建模。

\subsection{用户研究}

我们邀请了 24 名参与者（12 名数据分析师 $+$ 12 名前端开发者）进行 A/B 对比评估。每位参与者对来自 \ChartForge{} 和最佳基线（AMACE）的 10 组配对图表进行盲评，选择更满意的结果。

\begin{itemize}[leftmargin=*]
    \item \textbf{整体偏好}：\ChartForge{} \textbf{68.3\%} vs.\ AMACE 31.7\%（$p < 0.001$，二项检验）
    \item \textbf{数据分析师子组}：\ChartForge{} 71.2\% vs.\ AMACE 28.8\%（更偏好数据映射准确性）
    \item \textbf{前端开发者子组}：\ChartForge{} 65.4\% vs.\ AMACE 34.6\%（更偏好代码质量和可定制性）
\end{itemize}

%=====================================================================
\section{讨论}
%=====================================================================

\subsection{范式创新：从“生成代码”到“生成规约”}

\ChartForge{} 与现有 AIGC 图表方案的本质区别在于范式转变：\textbf{我们不对图表的具体实现代码（ECharts/Plotly/Vega-Lite 代码）建模，而是对一个平台无关的图表规约（Chart Specification）建模}。这与 React 的 Virtual DOM 哲学一脉相承——React 不直接操作浏览器 DOM，而是操作一个抽象的 Virtual DOM 树，由框架负责将抽象表示映射到具体平台。

这一范式的优势包括：

\begin{enumerate}[leftmargin=*]
    \item \textbf{平台无关性}：同一份 Chart Spec 可渲染为 Web（ECharts）、移动端（Swift Charts）、桌面端（Plotly）、AR 空间等多种形态。
    \item \textbf{可验证性}：Chart Spec 是结构化的形式化对象，可进行静态验证和定理证明，而生成的代码片段难以形式化验证。
    \item \textbf{生成效率}：规约空间远小于代码空间（规约约 200--500 tokens，代码约 2000--5000 tokens），大幅降低生成难度和延迟。
\end{enumerate}

\subsection{局限性}

\begin{enumerate}[leftmargin=*]
    \item \textbf{PCG 的图表类型覆盖}：当前 PCG 覆盖 12 类常见图表，对于高度定制化或新型图表（如网络图、3D 体渲染图）需要扩展产生式规则集。
    \item \textbf{SVAS 的主观性}：视觉美学的量化仍是一个开放问题，当前 SVAS 的视觉子分数依赖于启发式规则，未来可引入基于人类反馈的强化学习（RLHF）进行校准。
    \item \textbf{数据依赖}：PCG 的概率分布学习依赖于 ChartIntent-10K 数据集，对于特定垂直领域（如医学影像图表）需要领域适配。
\end{enumerate}

%=====================================================================
\section{结论}
%=====================================================================

本文提出了 \ChartForge{}——一个借鉴 React 声明式范式的 AIGC 图表控件生成框架。核心贡献包括：将图表生成从“代码生成”重构为“意图到规约的声明式映射”；设计了概率图表语法（PCG）为图表生成空间提供形式化基础；提出了语义-视觉对齐分数（SVAS）作为生成质量度量；构建了多阶段生成精炼管线（MS-GRP）实现迭代优化；定义了可组合图表代数（CCA）支撑声明式图表组合。

实验结果表明，\ChartForge{} 在 12 类图表控件上的生成准确率达 91.7\%，优于现有最佳方法 9.2 个百分点，同时生成延迟降低 69.2\%。未来工作将扩展 PCG 的图表类型覆盖、引入 RLHF 进行视觉偏好对齐、以及在 AR/VR 三维空间中实现沉浸式图表的 AIGC 生成。

%=====================================================================
% 参考文献
%=====================================================================

\begin{thebibliography}{99}

\bibitem{wongsuphasawat2016voyager}
K.~Wongsuphasawat, D.~Moritz, A.~Anand, \etal.
\newblock Voyager: Exploratory analysis via faceted browsing of visualization recommendations.
\newblock \textit{IEEE Transactions on Visualization and Computer Graphics}, 22(1):649--658, 2016.

\bibitem{walke2013react}
J.~Walke.
\newblock React: A {JavaScript} library for building user interfaces.
\newblock \textit{Facebook Engineering}, 2013.

\bibitem{tian2024chartgpt}
Y.~Tian, W.~Cui, D.~Deng, \etal.
\newblock {ChartGPT}: Leveraging {LLMs} to generate charts from abstract natural language.
\newblock \textit{IEEE Transactions on Visualization and Computer Graphics}, 2024.

\bibitem{koh2024c2}
W.~Koh, J.~H.~Yoon, \etal.
\newblock {C\textsuperscript{2}}: Scalable auto-feedback for {LLM}-based chart generation.
\newblock In \textit{Proceedings of NAACL}, 2025.

\bibitem{namgoong2025amace}
H.~Namgoong, J.~Jung, \etal.
\newblock {AMACE}: Automatic multi-agent chart evolution for iteratively tailored chart generation.
\newblock In \textit{Proceedings of EMNLP}, 2025.

\bibitem{satyanarayan2017vega}
A.~Satyanarayan, D.~Moritz, K.~Wongsuphasawat, J.~Heer.
\newblock {Vega-Lite}: A grammar of interactive graphics.
\newblock \textit{IEEE Transactions on Visualization and Computer Graphics}, 23(1):341--350, 2017.

\bibitem{wilkinson2005grammar}
L.~Wilkinson.
\newblock \textit{The Grammar of Graphics} (2nd ed.).
\newblock Springer, 2005.

\bibitem{vogelsang2019requirements}
A.~Vogelsang, M.~Borg.
\newblock Requirements engineering for machine learning: Perspectives from data scientists.
\newblock In \textit{Proceedings of RE Workshops}, 2019.

\bibitem{kim2020gemini}
Y.~Kim, J.~Heer.
\newblock {Gemini}: A grammar and recommender system for animated transitions in statistical graphics.
\newblock \textit{IEEE Transactions on Visualization and Computer Graphics}, 27(2):485--494, 2020.

\bibitem{moritz2019draco}
D.~Moritz, C.~Wang, G.~Nelson, \etal.
\newblock Formalizing visualization design knowledge as constraints: Actionable and extensible models in {Draco}.
\newblock \textit{IEEE Transactions on Visualization and Computer Graphics}, 25(1):438--448, 2019.

\bibitem{brown2020language}
T.~Brown, B.~Mann, N.~Ryder, \etal.
\newblock Language models are few-shot learners.
\newblock In \textit{Advances in Neural Information Processing Systems (NeurIPS)}, 2020.

\bibitem{devlin2019bert}
J.~Devlin, M.-W.~Chang, K.~Lee, K.~Toutanova.
\newblock {BERT}: Pre-training of deep bidirectional transformers for language understanding.
\newblock In \textit{Proceedings of NAACL}, 2019.

\bibitem{vaswani2017attention}
A.~Vaswani, N.~Shazeer, N.~Parmar, \etal.
\newblock Attention is all you need.
\newblock In \textit{Advances in Neural Information Processing Systems (NeurIPS)}, 2017.

\bibitem{sutton2018reinforcement}
R.~Sutton, A.~Barto.
\newblock \textit{Reinforcement Learning: An Introduction} (2nd ed.).
\newblock MIT Press, 2018.

\bibitem{maaten2008tsne}
L.~van~der~Maaten, G.~Hinton.
\newblock Visualizing data using {t-SNE}.
\newblock \textit{Journal of Machine Learning Research}, 9:2579--2605, 2008.

\end{thebibliography}

\end{document}
